\section{谁是龙头，谁只是跟风}

先分清角色，交易才不会跑偏。芒果超媒的价值在于它对应芒果TV平台敞口；它属于创业板，是当前观察题材热度的第一参照。主板名单里暂时找不到公司级项目关系股，所以也找不到合格的主板龙头。

\begin{exhibitbox}[展品：候选股分层]
\begin{tabularx}{\textwidth}{L{2.7cm}L{2.7cm}X L{2.2cm}}
\toprule
公司 & 角色 & 交易上代表什么 & 当前结论 \\
\midrule
芒果超媒（300413） & 平台敞口锚点 & 观察AI长剧叙事是否继续获得注意力。 & 创业板温度计。 \\
掌阅科技（603533） & 广义AI内容候选 & 观察主板能否承接AI内容扩散。 & 不能称为项目龙头。 \\
博纳影业（001330） & 影视扩散卫星 & 观察影视AI方向是否走强。 & 不承担龙头功能。 \\
中广天择（603721） & 内容/版权扩散卫星 & 观察内容数据方向是否走强。 & 不承担龙头功能。 \\
中文在线（300364） & 广义AI内容扩散 & 观察创业板AI内容强度。 & 不替代平台锚点。 \\
欢瑞世纪（000892） & 低纯度卫星 & 题材噪声可能更大。 & 排除主动观察池。 \\
\bottomrule
\end{tabularx}
\end{exhibitbox}

掌阅科技最容易被误读。它可以成为“AI内容”方向的主板容量股，却还未具备“AI长剧项目”主板核心的证据。两种叙事分开做，才不会把无法买到创业板龙头的焦虑，变成买错主板影子股。
